% 5 集合嵌套关系：N ⊂ Z ⊂ Q ⊂ R ⊂ C
\documentclass[tikz,border=4pt]{standalone}
\usepackage{ctex}
\usepackage{amsmath}
\usepackage{amssymb}
\usetikzlibrary{calc}
\begin{document}
\begin{tikzpicture}[thick]
  % 5 个嵌套矩形（左下到右上由内向外）
  \draw[fill=blue!8, draw=blue!60!black]   (-5.5,-3.5) rectangle (5.5,3.5);  % C
  \draw[fill=blue!16, draw=blue!60!black]  (-4.5,-2.9) rectangle (4.5,2.9);  % R
  \draw[fill=blue!28, draw=blue!60!black]  (-3.5,-2.3) rectangle (3.5,2.3);  % Q
  \draw[fill=blue!40, draw=blue!60!black]  (-2.5,-1.7) rectangle (2.5,1.7);  % Z
  \draw[fill=blue!55, draw=blue!60!black]  (-1.5,-1.1) rectangle (1.5,1.1);  % N

  % 集合符号标在每层左上角
  \node[anchor=north west, font=\large] at (-5.4, 3.4) {$\mathbb{C}$ 复数集};
  \node[anchor=north west, font=\large] at (-4.4, 2.8) {$\mathbb{R}$ 实数集};
  \node[anchor=north west, font=\large] at (-3.4, 2.2) {$\mathbb{Q}$ 有理数集};
  \node[anchor=north west, font=\large] at (-2.4, 1.6) {$\mathbb{Z}$ 整数集};
  \node[font=\large] at (0, 0.3) {$\mathbb{N}$};
  \node[font=\footnotesize] at (0, -0.3) {自然数集};

  % 嵌套关系标注（底部）
  \node[font=\small] at (0, -3.95)
    {$\mathbb{N} \subset \mathbb{Z} \subset \mathbb{Q} \subset \mathbb{R} \subset \mathbb{C}$};

  % 举例点（每层内一个代表元）
  \node[font=\scriptsize] at (-0.6, 0.85) {$0,1,2$};
  \node[font=\scriptsize] at (-1.9, 1.35) {$-3$};
  \node[font=\scriptsize] at (-2.9, 1.9) {$\tfrac{1}{2}$};
  \node[font=\scriptsize] at (3.5, 1.9) {$\sqrt{2},\pi$};
  \node[font=\scriptsize] at (4.5, 2.5) {$1+i$};
\end{tikzpicture}
\end{document}
